\section{Dữ liệu Healthcare}
 
\subsection{Các dạng dữ liệu}
 
Theo từ điển Cambridge, ``healthcare'' được định nghĩa là các hoạt động hoặc ngành nghề cung cấp dịch vụ y tế, trong tiếng Việt thường được hiểu là ``chăm sóc sức khỏe''. Trong bối cảnh khoa học dữ liệu, dữ liệu chăm sóc sức khỏe (Healthcare data) được định nghĩa là một hệ sinh thái dữ liệu đa chiều được hình thành từ sự tương tác liên tục giữa người bệnh và hệ thống y tế. Đây là tập hợp các thực thể thông tin phát sinh xuyên suốt hành trình của bệnh nhân, từ giai đoạn tầm soát, chẩn đoán, xét nghiệm cho đến lập kế hoạch điều trị. Nguồn dữ liệu này tồn tại dưới nhiều hình thái khác nhau bao gồm hồ sơ bệnh án điện tử (EHR), kết quả chẩn đoán hình ảnh, đơn thuốc, dữ liệu bảo hiểm và cả các tín hiệu sinh học từ các thiết bị giám sát sức khỏe (như IoT, EEG). Trong phạm vi khóa luận này, dữ liệu chăm sóc sức khỏe sẽ là nguồn tài nguyên số đóng vai trò đầu vào cho các mô hình dự đoán nhằm tối ưu hóa quyết định lâm sàng và nâng cao chất lượng sống cho cộng đồng.
 
Một số loại dữ liệu chăm sóc sức khỏe thường được thu thập bao gồm:
 
\paragraph{Hồ sơ sức khỏe điện tử (EHR) và hồ sơ bệnh án điện tử (EMR).}
Dữ liệu hồ sơ bệnh án điện tử (EMR) là phiên bản số hóa của hồ sơ bệnh án giấy của bệnh nhân tại phòng khám, gồm lịch sử điều trị cũng như y tế của một bệnh nhân trong phạm vi một cơ sở y tế. Trong khi đó, hồ sơ sức khỏe điện tử (EHR) lại được thiết kế để mở rộng và cung cấp cái nhìn tổng thể hơn về tình trạng sức khỏe của bệnh nhân. EHR chia sẻ thông tin với các nhà cung cấp, dịch vụ y tế và các tổ chức khác (như phòng khám đa khoa tư nhân, cơ sở chẩn đoán hình ảnh, nhà thuốc, cơ sở cấp cứu và bệnh viện) nhằm đảm bảo rằng thông tin sức khỏe của bệnh nhân có thể được truy cập và chia sẻ một cách an toàn giữa các môi trường chăm sóc sức khỏe khác nhau. Điều này tạo tính nhất quán, liền mạch cũng như chính xác trong quá trình trao đổi thông tin, đồng thời phục vụ cho mục đích chẩn đoán, xây dựng phác đồ điều trị và nâng cao chất lượng chăm sóc tổng thể.
 
\paragraph{Dữ liệu hành chính.}
Dữ liệu hành chính y tế là loại dữ liệu chủ yếu bao gồm thông tin tài chính và dữ liệu theo dõi hoạt động trong hệ thống chăm sóc sức khỏe, đóng vai trò quan trọng trong các chiến lược giám sát và quản lý. Loại dữ liệu này được thu thập chủ yếu nhằm phục vụ mục đích hành chính và thanh toán, đồng thời có thể được khai thác để phân tích việc cung cấp dịch vụ y tế, từ đó đánh giá ưu điểm, hạn chế cũng như hiệu quả chi phí của hệ thống. Ở một số quốc gia hoặc khu vực, dữ liệu hành chính còn bao gồm thông tin về xuất viện của bệnh nhân, được tổng hợp và báo cáo cho các cơ quan quản lý nhà nước hoặc tổ chức liên quan.
 
\paragraph{Dữ liệu yêu cầu thanh toán (Claims Data).}
Dữ liệu yêu cầu thanh toán là tập hợp các tương tác có thể lập hóa đơn giữa bệnh nhân có bảo hiểm và hệ thống cung cấp dịch vụ y tế, chủ yếu phục vụ mục đích chi trả hoặc bồi hoàn. Loại dữ liệu này bao gồm thông tin về chẩn đoán bệnh, các thủ thuật và xét nghiệm đã thực hiện, thời gian cung cấp dịch vụ, chi phí và địa điểm thực hiện. Thông thường, dữ liệu được chuẩn hóa thông qua các hệ thống mã như mã bệnh (ICD-10), mã thủ thuật (CPT) và mã định danh nhà cung cấp (NPI), nhằm đảm bảo tính nhất quán trong quy trình quản lý và thanh toán y tế. Những thông tin này góp phần phản ánh trải nghiệm thực tế của bệnh nhân cũng như quy trình cung cấp dịch vụ chăm sóc sức khỏe. Đồng thời, việc kết hợp khai thác dữ liệu yêu cầu thanh toán với các nguồn dữ liệu thực tế khác mở ra nhiều hướng tiếp cận mới trong nghiên cứu y tế.
 
\paragraph{Sổ đăng ký bệnh nhân/bệnh lý (Patient/Disease Registries).}
Sổ đăng ký là tập hợp các dữ liệu thứ cấp liên quan đến những bệnh nhân được chẩn đoán mắc một tình trạng, bệnh lý hoặc đã thực hiện một thủ thuật cụ thể. Đây là các hệ thống thông tin lâm sàng theo dõi những dữ liệu thiết yếu đối với các bệnh mạn tính như Alzheimer, ung thư, tiểu đường, bệnh tim và hen suyễn. Những nguồn thông tin này cung cấp các hiểu biết quan trọng trong việc quản lý và theo dõi tình trạng bệnh của bệnh nhân, đồng thời đóng vai trò quan trọng trong việc giám sát sau lưu hành của các loại dược phẩm. Một số ví dụ về sổ đăng ký bệnh lý bao gồm Mạng tương tác của Hiệp hội Alzheimer Toàn cầu (ALZConnected) hoặc Chương trình Quốc gia về đăng ký ung thư (NPCR).
 
\paragraph{Khảo sát sức khỏe (Health Surveys).}
Các dữ liệu khảo sát sức khỏe quốc gia được sử dụng để đánh giá tình trạng sức khỏe của dân số và ước tính tỷ lệ mắc bệnh. Loại dữ liệu này được thu thập nhằm phục vụ mục tiêu nghiên cứu và thường được công bố rộng rãi. Dữ liệu khảo sát sức khỏe đóng vai trò quan trọng trong quá trình ra quyết định khi xây dựng các chính sách và kế hoạch y tế, cung cấp thông tin chính xác về tình trạng dịch tễ học, mô hình sức khỏe, hành vi lối sống và trải nghiệm của bệnh nhân đối với các dịch vụ chăm sóc sức khỏe. Các nhà nghiên cứu y tế công cộng sử dụng dữ liệu khảo sát sức khỏe để phân tích hành vi liên quan đến sức khỏe và tình trạng tâm lý. Ngoài ra, loại dữ liệu này còn có giá trị trong việc xác định các tình trạng sức khỏe cụ thể hoặc các yếu tố nguy cơ trong cộng đồng, bao gồm việc sử dụng thuốc lá và rượu, chế độ ăn uống không lành mạnh và tình trạng ít vận động.
 
\paragraph{Dữ liệu thử nghiệm lâm sàng.}
Dữ liệu thử nghiệm lâm sàng được thu thập từ các nghiên cứu nhằm đánh giá các phương pháp điều trị y khoa, phẫu thuật hoặc hành vi trên bệnh nhân. Những dữ liệu này đóng vai trò là nguồn thông tin chính giúp khám phá các triệu chứng xuất hiện sớm, cách phòng ngừa, nâng cao chất lượng cuộc sống và cải thiện hiệu quả làm việc. Trong suốt quá trình thử nghiệm lâm sàng, nhiều loại dữ liệu khác nhau được thu thập và chuyển đổi thành các tập dữ liệu có thể phân tích nhằm giải quyết các câu hỏi nghiên cứu khác nhau. Sau đó, dữ liệu thử nghiệm lâm sàng có thể được sử dụng trong nhiều công bố khoa học và báo cáo phục vụ các nhóm đối tượng khác nhau.
 
\paragraph{Dữ liệu hệ gen.}
Dữ liệu hệ gen là dữ liệu bao gồm toàn bộ thông tin cấu trúc tế bào cần thiết cho hoạt động và sự phát triển từ bộ gen của sinh vật, bao gồm trình tự phân tử trong các gen, vai trò của từng gen, các yếu tố điều hòa kiểm soát biểu hiện gen và mối liên kết giữa các gen khác nhau. Thông tin hệ gen là nền tảng cho khoa học dữ liệu hệ gen, kết hợp giữa di truyền học và sinh học tính toán với phân tích dữ liệu thống kê và khoa học máy tính nhằm làm rõ cách dữ liệu di truyền ảnh hưởng đến sự phát triển và hoạt động của sinh vật, qua đó hỗ trợ nghiên cứu khoa học sự sống, chẩn đoán các bệnh di truyền và phát triển dược phẩm. Ví dụ, các nhà nghiên cứu sử dụng dữ liệu trình tự DNA để nghiên cứu bệnh tật và tìm ra các phương pháp điều trị cũng như thuốc mới.
 
\bigskip
\noindent\textbf{Dạng dữ liệu thường bắt gặp:}
 
\begin{itemize}[leftmargin=*, label={--}]
    \item \textbf{Dữ liệu cấu trúc:} Nhịp tim, huyết áp, cân nặng, kết quả xét nghiệm máu (Glucose, Creatinine, \ldots).
    \item \textbf{Dữ liệu phi cấu trúc (Unstructured Data):} Ghi chú của bác sĩ (clinical notes), bệnh sử, tóm tắt xuất viện. Dạng này thường cần kỹ thuật NLP (Xử lý ngôn ngữ tự nhiên) để khai thác.
    \item \textbf{Dữ liệu hình ảnh y khoa (Imaging Data):} Đây là các dữ liệu dạng pixel, đòi hỏi dung lượng lưu trữ lớn và kỹ thuật Deep Learning để phân tích. Ví dụ: X-quang, CT Scan, MRI, siêu âm, nội soi.
\end{itemize}
 
 
% ========== SECTION: Các vấn đề cần giải quyết ==========
\subsection{Các vấn đề cần được giải quyết trong lĩnh vực Healthcare Data}
 
\paragraph{Phân phối dữ liệu lệch và không đồng nhất (Skewed and Heteroskedastic Data).}
Dữ liệu chăm sóc sức khỏe thường có phân phối lệch rõ rệt, không tuân theo phân phối chuẩn đối xứng. Điều này xuất phát từ bản chất của các biến y tế như thời gian nằm viện, chi phí điều trị hay số lần tái khám --- vốn có xu hướng tập trung ở các giá trị thấp nhưng lại xuất hiện một số trường hợp cực đoan ở đầu phân phối cao. Ngoài ra, phương sai của sai số (heteroskedasticity) thường thay đổi theo giá trị của biến dự báo, vi phạm giả định cơ bản của nhiều mô hình hồi quy tuyến tính truyền thống. Nếu không được xử lý, các vấn đề này có thể dẫn đến ước lượng sai lệch, khoảng tin cậy không chính xác và khả năng tổng quát hóa của mô hình bị giảm sút đáng kể. Các kỹ thuật thường được áp dụng để khắc phục bao gồm biến đổi logarithm, biến đổi Box-Cox, hoặc sử dụng các mô hình có khả năng xử lý phân phối không chuẩn như hồi quy Poisson hay hồi quy Gamma.
 
\paragraph{Mối quan hệ phi tuyến (Nonlinear Relationships).}
Trong thực tế lâm sàng, mối quan hệ giữa các biến y tế hiếm khi tuân theo dạng tuyến tính đơn giản. Chẳng hạn, tác động của liều lượng thuốc lên hiệu quả điều trị, mối liên hệ giữa chỉ số BMI và nguy cơ mắc bệnh tim mạch, hay sự tương tác giữa tuổi tác và các chỉ số sinh hóa đều mang tính phi tuyến rõ rệt. Các mô hình tuyến tính truyền thống như hồi quy tuyến tính hay hồi quy logistic thường không đủ khả năng nắm bắt những mối quan hệ phức tạp này, dẫn đến sai lệch hệ thống trong dự đoán. Để giải quyết vấn đề này, các nhà nghiên cứu thường áp dụng các kỹ thuật như hồi quy đa thức, spline, mô hình cộng tổng quát hóa (GAM), hoặc các phương pháp học máy có khả năng học biểu diễn phi tuyến như cây quyết định, random forest và mạng nơ-ron nhân tạo.
 
\paragraph{Sự mất cân bằng giữa các lớp trong phân loại (Imbalanced Classes).}
Mất cân bằng lớp là một thách thức đặc trưng và phổ biến trong dữ liệu y tế, xuất phát từ sự chênh lệch tự nhiên giữa tỷ lệ mắc bệnh và tỷ lệ không mắc bệnh trong dân số. Trong nhiều bài toán phân loại lâm sàng, lớp dương tính (ví dụ: bệnh nhân mắc bệnh hiếm gặp, ca tử vong, hoặc biến cố bất lợi) thường chiếm tỷ lệ rất nhỏ so với lớp âm tính. Điều này khiến các mô hình học máy có xu hướng thiên vị về phía lớp đa số, đạt độ chính xác tổng thể cao nhưng lại thất bại trong việc nhận diện đúng các trường hợp thuộc lớp thiểu số --- vốn là mục tiêu quan trọng nhất trong bối cảnh y tế. Theo nghiên cứu của Tyagi \& cộng sự (2022), một khung phương pháp hiệu quả để giải quyết vấn đề này bao gồm việc lựa chọn mô hình nền phù hợp (Logistic Regression, SVM, Decision Tree) trước, sau đó áp dụng các kỹ thuật tái cân bằng dữ liệu như SMOTE với tham số alpha được tối ưu hóa, Random Oversampling/Undersampling, hay điều chỉnh ngưỡng phân loại thông qua học nhạy trọng số.
 
\paragraph{Dữ liệu bị thiếu và hiện tượng bỏ cuộc (Missing Data and Dropout).}
Dữ liệu bị thiếu là vấn đề gần như không thể tránh khỏi trong các nghiên cứu y tế, đặc biệt trong các bộ dữ liệu thu thập từ hồ sơ bệnh án điện tử hoặc các thử nghiệm lâm sàng dài hạn. Dữ liệu có thể bị thiếu theo nhiều cơ chế khác nhau: hoàn toàn ngẫu nhiên (MCAR --- \textit{Missing Completely At Random}), ngẫu nhiên có điều kiện (MAR --- \textit{Missing At Random}) hoặc không ngẫu nhiên (MNAR --- \textit{Missing Not At Random}), và mỗi cơ chế đòi hỏi chiến lược xử lý khác nhau. Bên cạnh đó, hiện tượng bỏ cuộc (dropout) trong các nghiên cứu theo dõi dọc --- khi bệnh nhân ngừng tham gia nghiên cứu vì nhiều lý do như tử vong, di chuyển hoặc từ chối tiếp tục --- có thể tạo ra thiên kiến lựa chọn nghiêm trọng nếu không được kiểm soát. Ngoài phương pháp điền thế đơn giản như mean/median imputation, các kỹ thuật nâng cao hơn như điền thế đa biến (Multiple Imputation by Chained Equations --- MICE), mô hình hóa dữ liệu bị thiếu, hay điều chỉnh thiên kiến lựa chọn trong các thử nghiệm mục tiêu mô phỏng (như đề xuất của Fernández-García \& cộng sự, 2024) đang ngày càng được ứng dụng rộng rãi.
 
\paragraph{Vấn đề ``Lời nguyền đa chiều'' trong dữ liệu lớn (High-Dimensional Data).}
Với sự phát triển của công nghệ giải trình tự gen thế hệ mới, dữ liệu hệ gen trong y tế hiện nay có thể chứa hàng triệu biến số (SNP, biểu hiện gen) trong khi số lượng mẫu bệnh nhân thường bị giới hạn đáng kể. Hiện tượng này, được gọi là ``lời nguyền đa chiều'' (\textit{curse of dimensionality}), khiến không gian dữ liệu trở nên cực kỳ thưa thớt, các phép đo khoảng cách mất ý nghĩa thống kê, và nguy cơ overfitting tăng cao đột biến. Hệ quả là nhiều mô hình học máy truyền thống không thể hoạt động hiệu quả khi số chiều dữ liệu vượt xa số quan sát. Các giải pháp phổ biến bao gồm kỹ thuật giảm chiều dữ liệu như PCA (Principal Component Analysis), t-SNE, UMAP, cùng với các phương pháp chọn lọc đặc trưng (feature selection) và chính quy hóa (regularization) như LASSO và Ridge Regression nhằm loại bỏ các biến không có giá trị dự báo và tập trung vào các đặc trưng thực sự có ý nghĩa lâm sàng.
